%% ma: L12-B16-0006
%% lop: 12
%% chuong: 5
%% bai: 16
%% dang: 12.16.4
%% loai: TN4
%% muc: TH
%% dap_an: "C"
%% nguon: "(NBV)-ĐỀ SỐ 18-MỖI NGÀY 1 ĐỀ THI 2025.pdf (D18, MỖI NGÀY 1 ĐỀ - Đề số 18), Phần I Câu 10"
%% kien_thuc: "Bài 16 (góc giữa đường thẳng và mặt phẳng), Bài 8 (tọa độ vectơ)"
%% trang_thai: cho-duyet
%% ly_do: "Dạng mới đề xuất 12.16.4 (côsin góc giữa đường thẳng đi qua hai điểm và mặt phẳng tọa độ), chờ giáo viên duyệt."
%% ghi_chu: "Phương án C của đề gốc in là căn 337 chia 29 (lỗi in); giải độc lập cho cos a = căn(13/29) = căn 377 chia 29 (≈ 0,6695), khớp đáp án gốc C nên bản kho sửa phương án C thành căn 377 chia 29."
\begin{ex}
Trong không gian $Oxyz$, cho hai điểm $A(1;3;2)$ và $B(4;5;6)$. Gọi $\alpha$ là góc giữa đường thẳng $AB$ và mặt phẳng $(Oxy)$. Giá trị của $\cos\alpha$ bằng
\choice{$\dfrac{4\sqrt{29}}{29}$}{$\dfrac{16}{29}$}{\True $\dfrac{\sqrt{377}}{29}$}{$\dfrac{13}{29}$}
\loigiai{Ta có $\overrightarrow{AB}=(3;2;4)$, $|\overrightarrow{AB}|=\sqrt{29}$. Mặt phẳng $(Oxy)$ có vectơ pháp tuyến $\vec k=(0;0;1)$ nên $\sin\alpha=\dfrac{|\overrightarrow{AB}\cdot\vec k|}{|\overrightarrow{AB}|\cdot|\vec k|}=\dfrac4{\sqrt{29}}$. Vì $0^\circ\le\alpha\le90^\circ$ nên $\cos\alpha=\sqrt{1-\dfrac{16}{29}}=\sqrt{\dfrac{13}{29}}=\dfrac{\sqrt{377}}{29}$. Chọn C.}
\end{ex}
